% !TEX root = ../AGAO.tex

\section{Classical Algorithms for Maximal Flow 2}
We have the same setup as before, but we will approach it a bit differently.

\subsection{Blocking Flows}

Let $\ff$ be any feasible flow in $G$ and $G_f$ its residual graph. 

We can partition all the vertices of $G_f$ by distance from $s$, so define the \textbf{level} of a vertex $v$ to be the length of the shortest path from $s$ to $v$. Denote that by $\ell (v)$ or $\ell_{G_f} (v)$, depending on the context. We observe that an edge can only bring you up one layer, but can drop you down multiple layers.

Here we will try to focus on \textit{short} paths. 

\begin{definicija}
    We call an edge \emph{admissible} if it brings us up exactly one level (edge $(u,v)$, $\ell(u) + 1 = \ell(v)$.

    We take all the admissible edges of $G_f$ and the graph induced by them is calleed the \emph{level graph} of $G_f$.

    A \emph{blocking flow} in $G_f$ is a feasible flow $\Tilde{\ff}$ that uses only admissible edges while also saturating at least one edge along any $s-t$ path in the level graph of $G_f$.
    
\end{definicija}

It is blcking because of the last condition. Basically it blocks every shortest path. Of course it might not be maximum flow in a graph.


\subsection{Dinic's Algo}

This will be the essence of Dinic's Algorithm for calculating max flow. We start with $\ff = 0$ and repeatedly add blocking flows until no more $s-t$ paths exist in $G_f$.

Since it is a special case of Ford-Fulkerson we know that we have the same properties.

\begin{lema}
    Let $\ff$ be a feasible flow and $\Tilde{\ff}$ a blocking flow in $G_f$. If we define $\ff ' = \ff + \Tilde{\ff}$, then 
    \begin{align*}
        \ell_{G_{\ff'}}(t) \geq \ell_{G_\ff}(t) +1
    \end{align*}
\end{lema}

\obvs

From this we immediately get the following

\begin{izrek}
    Dinic's algo terminates in $O(n)$ iterations
\end{izrek}

\obvs

For graphs with unit capacities we can actually get better bounds:

\begin{izrek}
    On unit capacity graphs, Dinic's algo terminates in
    \begin{align*}
        O ( \min (m^{1/2}, n^{2/3}))
    \end{align*}
\end{izrek}

\begin{proof}
    \begin{enumerate}
        \item Suppose we run Dinic's algo for $k$ iteration. Now take $\Tilde{\ff}$ to be the optimal flow in $G_f$ (watch out, not in $G$) and since every path from $s$ to $t$ now takes at least $k$ edges and routes exactly $1$ unit of demand, we are using total edge capacity of 
        atleast $k \cdot \text{val} (\Tilde{\ff})$. But since the total number of edges (with nonzero capacity) in $G_f$ is still the same as in $G$ 
        the total capacity of $G_f$ is at most $m$, so 
        $\textbf{val} (\tff) \leq \frac{m}{k}$

        Recalling FF algorithms, this implies we wil terminate after atmost another $\frac m k$ iterations. Now we want to minimize $k + \frac m k$ and we see that this is optimized when $k = \sqrt{m}$ and we get $O(\sqrt{m})$ steps.

        \item Again suppose we run it for $k$ iterations to get a flow $\ff$. Graph $G_f$ partitions the vertices into levels, so the sink must be at the level atleast $k$, so every layer from $1$ to $k$ has to be nonempty. And in $G_f$ we are only interested in the shortest paths to $t$, so we can forget about any vertices that are further away than $t$.

        We see that we can't have more than $\frac k 2$ levels with size $|D_i| \geq \frac{2n}{k}$, as we have atmost $n-1$ elements in there, so there are more than $\frac k 2$ levels with $|D_i| < \frac{2n}{k}$, so there are atleast two consecutive levels with that and thus there are atmost $\frac{4n^2}{k^2}$ edges going there and from min cut we get that the remaining val is atmost that, so this also bounds the number of iterations left over.

        So we need to do atmost $k + \frac{4n^2}{k^2}$ iterations, which is minimized at $k = 2n^{2/3}$ and the result is $O(n^{2/3}$.
    \end{enumerate}
    So we have two bounds and we can take the lower one.
\end{proof}


\subsection{How to find Blocking Flows}

So we havent talkd about how to find Nemo yet.

But we can do that using repeated DFS in the level graph.

Though we need to have the level graph first, which we can achieve with BFS I guess and classify everything with that. And I believe that should take $O(m+n)$, which will be negligible at the end.

Once we have this we just go deeper to try to hit $t$. If we arrive to the correct level without hitting $t$, then we switch to the new path and delete the first edge that we tried. Now the complexity of this should be $O(mn)$

The algorithms is decomposed into two algorithms

\begin{enumerate}
    \item FINDBLOCKINGFLOW
    \item DFS(u, H,t)
\end{enumerate}



\begin{lema}
    For general graphs, FINDBLOCKINGFLOW returns a blocking flow in $O(mn)$ time. Hense Dinic's algo runs in $O(n^2m)$ time on general capacity graphs.
\end{lema}

\begin{proof}
    Each useless edge will get added and deleted only once, so we get rid of them in $O(m)$ time, since there are at most that many useless edges.

    Now we look at the useful ones. When we find a path, we will have at most $n-1$ edges, since that is the longest path, so that will be $O(n)$. You might think, that we need to be careful because we do things with the good edge even if we go to a useless one, but that is already taken care by the general $O(m)$ before. And we have at most $m$ different paths here, since every path saturates at least one edge. So we have $O(m+n) + O(m) + O(mn) = O(mn)$

    And because the path can't be longer than $n$, we will repeat all the steps at most $n$ times with different level graphs, so that brings us to $O(n^2m)$

    We note that the augmented paths can still interesect at some paths, but each will have atleast one saturated edge.
\end{proof}

\begin{lema}
    For unit capacity graphs, FINDBLOCKINGFLOW returns a blocking flow in $O(m)$ time and Dinic Algo runs in $O(\min (m{3/2}, mn^{2/3})$ time.
\end{lema}

\begin{proof}
    Here we know that if a path goes through an edge, no other path can go through it, so each edge will be traversed atmost once by DFS.
\end{proof}

There is a result on bipartite matching graphs where Dinic's algo runs in $O(|E| \sqrt{|V|})$, but we will not look at it.


\subsection{Min Cut is a Linear Program}

Remember the min cut as taking a subset $S$ of $V$ and its complement in $V$ and looking at edges starting from $S$ going to the complement. We want to find a set that minimizes the capacity of those edges.

This should be equivalent to the following minimization program:

\begin{equation}
    \begin{aligned}
    \min _{\mathbf{x}\in \R^V} &  \quad \sum_{e\in E} \ccc (e) \max(b_e ^\top \mathbf{x}, 0)\\
    \text{with} & \quad  \mathbf{x}(s) = 0\\
    & \quad \mathbf{x}(t) = 1 \\
    & \quad 0 \leq \mathbf{x} \leq 1
\end{aligned}
\end{equation}

Where $b_e^\top x = x(v) - x(u)$ for $e = (u,v)$

This might not look like a proper linear program because of the $\max$ inside there, but it actually is if we introduce another variable $\mathbf{u} \in \R^E$

\begin{equation}
    \begin{aligned}
    \min _{\mathbf{x}\in \R^V, \mathbf{u} in \R^E} &  \quad c^\top \mathbf{u}\\
    \text{with} & \quad  \mathbf{x}(s) = 0\\
    & \quad \mathbf{x}(t) = 1 \\
    & \quad 0 \leq \mathbf{x} \leq 1\\
    & \quad \mathbf{u} \geq 0\\
    & \quad \mathbf{u} \geq \mathbf{B^\top x}
\end{aligned}
\end{equation}

Somehow we should see that we don't neccesarily need to bound x between $0$ and $1$, but I don't really care rn.

There is a proof of that in the script.